\section{Limitations}
\label{sec:limitations}

% JLEO-R1: rewrite limitations around overlap instability, holdout weakness, and narrow external validity
We catalog the limitations of the design under three headings:
identification gaps that the institutional design and supporting
diagnostics narrow but do not close, scope constraints inherent to
the BEC setting, and operational caveats for deployment. Each
limitation pairs the residual threat with the bound we report and
the next step that would relax it.

\paragraph{Within-market selection.}
The baseline FL--price coefficient is not robust to stricter
observable overlap. The main within-item estimate is
$\valMatchBaselineCoef$, but exact overlap on item-group, year,
modality, PBU size, and tender-value quartile turns it to
$\valMatchOverlapCoef$; adding reference-price bins leaves it at
$\valMatchOverlapRefCoef$; ATT weighting on observables drives it
to $\valMatchPSCoef$. This pattern is stronger than a generic
caveat about omitted variables. It says directly that the
estimated sign depends on which parts of the sample supply the
comparison. Cinelli--Hazlett still bounds how strong an
unobserved confounder would need to be to erase the baseline,
but the rereview version treats the price result as a
conditional association only. \emph{Next step:} a design with
sharper exogenous treatment timing or richer pre-treatment
market covariates.

\paragraph{Pre-trends in the FL-entry event study.}
Pre-event coefficients at $t = -2$ and $t = -3$ are positive and
similar in magnitude to the post-event coefficient
(Figure~\ref{fig:event_study}, \ref{app:did}). Two readings are
observationally equivalent: cartels select into already-high-price
markets and then deploy cover bidders, or unobserved
within-market dynamics predate FL arrival. The paper's
diagnostic positioning does not require choosing between them;
top-5-grade causal identification would.
\citet{rambachan2023more} sensitivity bounds the maximum
pre-trend slope under which conclusions hold; we report the
bounds as failing the conventional benchmarks for the
event-study estimates, which is why we report the staggered DiD
as a transparency artifact rather than as a robustness check
(Section~\ref{sec:robustness}). \emph{Next step:} a sharply
exogenous treatment timing the FL setting does not provide.

\paragraph{FL persistence as environment marker.}
The pooled-sample cutoff is not stable across windows. Training
the threshold only on 2009--2016 yields a cutoff of
$\valThresholdTrain$, below the full-sample statistical cutoff
of $\valThresholdStat$, and reduces the firm-level AUC to
$\valAUCStrictFirmFL$. On 2017--2019 items the frozen binary
rule drops to $\valAUCStrictItemFL$, while the continuous score
retains $\valAUCStrictItemTC$. FL is therefore properly read as
an environment-level marker and binary trigger, not as a stable
firm type or a time-invariant structural classification.
\emph{Next step:} periodic retraining of the cutoff or direct
deployment of the continuous ranking score.

\paragraph{Zero-win threshold dependence.}
Relaxing the zero-win condition to include firms with 1--2\%
win rates substantially attenuates the price coefficient
(Table~\ref{tab:fl_lowwinrate}, \ref{app:online_robustness}).
The bright line at zero is doing real work: a sophisticated
cartel could evade the screen by allowing cover bidders the
occasional win. \emph{Next step:} adversarial-aware screen
design---weighted scores rather than indicator thresholds, or
periodic threshold updating---which the
\citet{harrington2008detecting} cat-and-mouse logic recommends
for any proactive enforcement tool.

\paragraph{CADE-overlap conditional indistinguishability.}
Conditional on participation intensity, FL firms are
statistically indistinguishable from non-FL always-losers in
CADE-proximity rate ($p = 0.93$;
Section~\ref{sec:ground_truth}). The unconditional
$\valMultThreeFive\times$ excess we report
(Section~\ref{sec:cade}) is therefore driven by participation
volume, not by the FL classification per se. This is a feature
of the screen, not a bug---it flags firms that participate
heavily in suspicious markets without requiring the analyst to
distinguish participation-driven proximity from direct cartel
involvement---but it means the CADE check validates
\emph{market-level exposure}, not firm-level cartel membership.
\emph{Next step:} matching FL flags to leniency-driven
firm-level identifications when the latter become available.

\paragraph{Prospective ground truth.}
Eight of the twelve CADE cases were adjudicated 2020--2025,
after the BEC sample window closes. A strictly contemporaneous
validation would restrict ground truth to the four pre-2020
cases, leaving too few positives for inference. The
prospective framing we adopt is methodologically defensible
(future enforcement cannot leak into pre-2019 participation
patterns) but mixes firms whose collusive behavior was active
during the BEC window with firms whose cases were filed
earlier and adjudicated later. \emph{Next step:} validation
against contemporaneous-window adjudications as more cases
clear adjudication; in the meantime, the strict temporal
holdout in Table~\ref{tab:strict_train_threshold} is the most
honest deployment benchmark we can provide.

\paragraph{Structural model is a calibration.}
The mixture-model regime selection and implied markup
(\ref{app:structural_id}, Online) calibrate cover-bid
distributions but do not identify the markup or the causal
price effect independently of the reduced-form coefficients.
Welfare quantification in \ref{app:extensions} (Online)
mechanically applies price coefficients to FL-present spending
and \emph{assumes} causality; we report the resulting
$\valWelfareLowPct$--$\valWelfareHighPct\%$-of-spending range as illustrative, not as a
policy estimate.

\paragraph{External validity.}
% JLEO-R2: scope statement strengthened with falsifiable prediction
The paper's results come from a single Brazilian platform with a
very specific data structure and product mix. In BEC, roughly
\valExtCommodityShare\ of items are commodity-like and only
\valExtServiceShare\ are service-like under a transparent text
classification. Detection is also modality-skewed:
pregão-primary always-losers yield AUC $\valExtPregAUC$ versus
$\valExtConvAUC$ for convite-primary firms. These facts do not
invalidate the screen, but they narrow its natural external
domain.

We make this scope claim falsifiable by stating where the
construct's empirical content is most likely portable and where
it is not. The construct is most likely portable to procurement
environments sharing three structural features: (i) commodity-heavy
item composition (allowing repeated participation by the same
firms across many tender-items), (ii) electronic platforms with
publicly auditable participant lists, and (iii) institutional
asymmetry between winner and loser sides of the bidding pool
(e.g., minimum-bidder rules in convite-style modalities, or
analogous quorum requirements in other jurisdictions). The
construct is least likely portable to (i) heterogeneous works
contracts where each tender involves different firms, (ii)
private-sector procurement where participant lists are not
disclosed, and (iii) jurisdictions where bid microdata are
routinely audited and bid-distribution screens are operational
\citep{imhof2019detecting,wallimann2023machine}. The construct's
value is in cheaply screening procurement environments where
bid microdata are unavailable---roughly the OECD-procurement-without-microdata
frontier surveyed by \citet{coviello2017tenure} and the
participation-disclosure regimes documented by
\citet{decarolis2020rules}. Generalization to ComprasNet, other
Brazilian states, or OECD platforms requires external
replication with the same scope discipline.

\begin{table}[htbp]
\centering
\caption{External Validity Scope: Modality and Commodity-Service Boundaries}
\label{tab:external_validity_scope}
\begin{threeparttable}
\small
\begin{tabular}{lllcc}
\toprule
Dimension & Group & Metric & Value & 95\% CI \\
\midrule
coverage & Commodity & share\_items & $0.887$ & --- \\
coverage & Service & share\_items & $0.113$ & --- \\
dominant\_item\_class & Commodity & auc\_log\_tc & $0.905$ & $[0.879, 0.931]$ \\
dominant\_item\_class & Service & auc\_log\_tc & $0.966$ & $[0.961, 0.971]$ \\
item\_class\_price & Commodity & coef\_losers & $0.045$ & $[0.002, 0.088]$ \\
item\_class\_price & Service & coef\_losers & $0.128$ & $[0.024, 0.231]$ \\
modal\_primary\_auc & convite\_primary & auc\_log\_tc & $0.816$ & $[0.758, 0.874]$ \\
modal\_primary\_auc & pregao\_primary & auc\_log\_tc & $0.952$ & $[0.946, 0.958]$ \\
\bottomrule
\end{tabular}
\begin{tablenotes}
\small
\item \textit{Notes:} Item classes use a transparent text rule on BEC item descriptions. Rows labelled \texttt{modal\_primary\_auc} reuse the always-loser firm split by primary modality; rows labelled \texttt{item\_class\_price} report the within-item price association separately for commodity-like and service-like items. The point is not to claim exportability to all procurement, but to document that the paper speaks most naturally to standardized-goods environments and selected simple services.
\item \textit{Source:} \texttt{scripts/52\_external\_validity\_scope.R}.
\end{tablenotes}
\end{threeparttable}
\end{table}

